\clearpage
% \setcounter{page}{1}
% \maketitlesupplementary
\appendix

\section{Implementation Details}
\label{app:exp-settings}

\begin{table}[t]
\small
\centering
\begin{adjustbox}{max width=\linewidth}
\begin{tabular}{y{60}x{40}x{28}x{28}x{33}}
\toprule
configs & iMF-B & iMF-M & iMF-L & iMF-XL \\
\midrule
params (M)
            & 89  & 174 & 409  & 610 \\
depth
            & 12  & 24  & 32   & 48 \\
hidden dim
            & 768 & 768 & 1024 & 1024 \\
attn heads
            & 12  & 12  & 16   & 16 \\
patch size
            & \multicolumn{4}{c}{$2{\times}2$} \\
aux-head depth  
            & \multicolumn{4}{c}{8} \\
class tokens 
            & \multicolumn{4}{c}{8} \\
time tokens 
            & \multicolumn{4}{c}{4} \\
guidance tokens  
            & \multicolumn{4}{c}{4} \\
interval tokens  
            & \multicolumn{4}{c}{4} \\
linear layer init
            & \multicolumn{4}{c}{ $\mathcal{N}(0, \sigma^2)$, $\sigma^2=0.1 / \mathtt{fan\_in}$ } \\
\midrule
epochs
            & 240$^\dagger$ / 640 & 640 & 640 & 800    \\
batch size  
            & \multicolumn{4}{c}{256$^\dagger$ / 1024} \\
learning rate 
            & \multicolumn{4}{c}{0.0001} \\
lr schedule
            & \multicolumn{4}{c}{constant} \\
lr warmup~\cite{largesgd} 
            & \multicolumn{4}{c}{10 epochs} \\
optimizer
            & \multicolumn{4}{c}{Adam \cite{adam}} \\
Adam $(\beta_1, \beta_2)$
             & \multicolumn{4}{c}{(0.9, 0.95)} \\
weight decay 
            & \multicolumn{4}{c}{0.0} \\
dropout
            & \multicolumn{4}{c}{0.0} \\
ema decay 
            & \multicolumn{4}{c}{0.9999} \\
\midrule
ratio of $r{\neq}t$ 
            & \multicolumn{4}{c}{50\%} \\
$(t,r)$ cond 
            & \multicolumn{4}{c}{$t-r$} \\
$t,r$ sampler 
            & \multicolumn{4}{c}{logit-normal($-0.4$, $1.0$)} \\
\midrule
cls drop \cite{cfg} 
            & \multicolumn{4}{c}{0.1} \\
CFG dist $\beta$ 
            & 1 & 2 & 2 & 2 \\ 
\bottomrule
\end{tabular}
\end{adjustbox}
\caption{\textbf{Configurations and hyper-parameters}. $^\dagger$: these are for ablation studies.}
\label{tab:imagenet-configs}
\end{table}


The configurations and hyper-parameters are summarized in \cref{tab:imagenet-configs}. Our implementation is based on the public codebase of original MF, which is based on JAX and TPUs.\footnotemark{}

\footnotetext{\url{https://github.com/Gsunshine/meanflow}}

\paragraph{Auxiliary head for $v_\theta$.} In \cref{sec:v-pred}, we have introduced an auxiliary head for modeling $v_\theta$, which produces the input to the $\mathtt{JVP}$ computation. This auxiliary head shares most layers and computation with the main network $u_\theta$ and only differs in the last $L$ layers (we set $L=8$ in all cases). The output of this auxiliary head is to predict the marginal $v$.

Unlike the variant using the boundary condition of $u$ (that is, $v_\theta(z_t)\,=u_\theta(z_t, t, t)$), the unshared layers in the auxiliary head receive no gradient if not attached to any loss. To address this issue, we append an auxiliary loss $\|v_\theta - (e-x)\|^2$ to this head, which is the Flow Matching loss. The output of this head is only for the $\mathtt{JVP}$ computation and is not used at inference time.
As adaptive weighting~\cite{mf} is used in the MF loss, we also apply it to this auxiliary loss.

\paragraph{CFG conditioning.} In \cref{sec:method-cfg}, we have discussed CFG using the standard form \cite{cfg}: $v_{\text{cfg}} = \omega\,v(z_t\mid\mathbf{c}) + (1-\omega)\,v(z_t \mid \varnothing)$ (\cref{eq:cfg}).
The original MF paper \cite{mf} derives a relationship between $v_{\text{cfg}}$ and its resulting average velocity field, which can be simplified as (see Eq.~(21) in \cite{mf}):
\begin{equation}
    v_{\text{tgt}}=(e-x)+(1-\frac{1}{\omega})\big(u^\text{cfg}_\theta(z_t\mid t,t,\mathbf{c}) - u^\text{cfg}_\theta(z_t\mid t,t, \varnothing)\big),
\end{equation}
where ``$\varnothing$'' is to emphasize the unconditional field. 
Here, $\omega$ is the ``effective guidance scale'' \cite{mf}, which plays the same role as in standard CFG.
Specifically, when $\omega=1$, this equation degenerates to the no-CFG case. 
We adopt this formulation.
See the pseudocode in \cref{alg:cfg_code}. 

\definecolor{codeblue}{rgb}{0.25,0.5,0.5}
\definecolor{codesign}{RGB}{0, 0, 255}
\definecolor{codefunc}{rgb}{0.85, 0.18, 0.50}


\lstdefinelanguage{PythonFuncColor}{
  language=Python,
  keywordstyle=\color{blue}\bfseries,
  commentstyle=\color{codeblue},
  stringstyle=\color{orange},
  showstringspaces=false,
  basicstyle=\ttfamily\small,
  literate=
    {+}{{\color{codesign}+ }}{1}
    {-}{{\color{codesign}- }}{1}
    {*}{{\color{codesign}* }}{1}
    {/}{{\color{codesign}/ }}{1}
    {sample_t_r_cfg}{{\color{codefunc}sample\_t\_r\_cfg}}{1}
    {randn_like}{{\color{codefunc}randn\_like}}{1}
    {jvp}{{\color{codefunc}jvp}}{1}
    {stopgrad}{{\color{codefunc}stopgrad}}{1}
    {metric}{{\color{codefunc}metric}}{1}
    % {None}{{\color{codeblue}None}}{0}
}

\lstset{
  language=PythonFuncColor,
  backgroundcolor=\color{white},
  basicstyle=\fontsize{8pt}{8.8pt}\ttfamily\selectfont,
  columns=fullflexible,
  breaklines=true,
  captionpos=b,
}

\begin{algorithm}[t]
\centering
\caption{{improved MeanFlow}: training guidance.\\
{\scriptsize Note: in PyTorch and JAX, \texttt{jvp} returns the function output and JVP.}}

\begin{minipage}{0.99\linewidth}
\begin{lstlisting}[language=PythonFuncColor, escapechar=`]
# fn(z, t, r, w, c): function to predict u
# x: training batch
# c: condition batch

t, r, w = sample_t_r_cfg()
e = randn_like(x)

z = (1 - t) * x + t * e

# cls cond and cls uncond v
v_c = fn(z, t, t, w, c)
v_u = fn(z, t, t, w, None)

# Compute CFG target (same as orig MF)
v_tgt = (e - x) + (1 - 1 / w) * (v_c - v_u)

# Use predicted v_c to compute dudt
u, dudt = jvp(fn,  (z, r, t, w, c),  
                        (v_c, 0, 1, 0, 0))

# Compute compound function V
V = u + (t - r) * stopgrad(dudt)
error = V - stopgrad(v_tgt)

loss = metric(error)
\end{lstlisting}
\end{minipage}
\label{alg:cfg_code}
\end{algorithm}


When using CFG-conditioning, we need to randomly sample the scale $\omega$ for each training sample.
First, we set a sampling range of $\omega$:  $[1.0,\, \omega_{\max}]$, where we fix $\omega_{\max}=8.0$.
Note that $\omega=1$ degenerates to the no-CFG case.
Then we sample $\omega$ from a power distribution that biases towards smaller $\omega$ values: $\omega \sim p(\omega)\propto\omega^{-\beta}$,
where $\beta$ controls the skewness (we use $\beta = 1$ or $2$, \cref{tab:imagenet-configs}).

When using CFG-conditioning to support guidance interval \cite{interval}, we randomly sample $t_{\min}$ and $t_{\max}$ from $\mathcal{U}[0,0.5]$ and $\mathcal{U}[0.5,1.0]$ respectively, where $\mathcal{U}$ is the uniform distribution. During training, when $t$ falls outside of $[t_{\min}, t_{\max}]$, CFG is turned off by setting $\omega=1$.
The set of sampled values of $\Omega=\{\omega, t_{\min}, t_{\max}\}$ is provided to the network as extra conditioning.

\paragraph{In-context conditioning.} Our models are conditioned on time steps $r, t$, class $\mathbf{c}$, and CFG factors $\Omega=\{\omega, t_{\min}, t_{\max}\}$. All continuous-valued conditions (\eg $\omega$) are processed by standard positional embedding \cite{transformer}, similar to how $t$-conditioning is handled in continuous-time diffusion/flow models.
Each type of these conditions is processed by a 2-layer MLP. All conditions are replicated into multiple tokens: the number of replications for each type is in \cref{tab:imagenet-configs}. All replicated tokens are added with learnable embeddings indicating their types of conditions (analogous to position embedding along the sequence), and then are concatenated along the sequence axis with the image tokens (see \cref{fig:network}).

\paragraph{Removing adaLN-zero.} When using in-context conditioning, our model removes the standard adaLN-zero \cite{dit} that is parameter-heavy.
We adopt the zero residual-block initialization \cite{largesgd}, which adaLN-zero \cite{dit} also follows. Specifically, for any residual block \cite{resnet} in a Transformer \cite{transformer} with the form of $x+F(x)$, the last operation in $F(x)$ is always a learnable per-channel scale $\gamma$, where $\gamma$ is initialized as zero. As such, the initial state of a residual block is always identity mapping. The initialization of adaLN-zero \cite{dit} was based on the same principle.

For all other linear projection layers in the Transformer blocks, we use a Gaussian initialization $\mathcal{N}(0, \sigma^2)$, where $\sigma^2=0.1 / \mathtt{fan\_in}$ and $\mathtt{fan\_in}$ is the input channel number. This can be implemented in popular libraries by $\sigma=\mathtt{gain} / \sqrt{\mathtt{fan\_in}}$ where $\mathtt{gain}=\sqrt{0.1}\approx 0.32$. This initialized $\sigma$ is more conservative than common gain-controlled initializations, where $\mathtt{gain}$ is 1 or $\sqrt{2}$. In our preliminary experiments, this initialization converges faster when our block becomes different from the adaLN-zero block.

\paragraph{Evaluation}.
We sample 50,000 samples and compute FID against the ImageNet training set (\ie FID-50K). We sample 50 images per class for the FID evaluation. For each of our models where CFG-conditioning is enabled, we report the FID results using the optimal guidance scale and interval.

\section{Additional Qualitative Results} 
\label{app:vis}
We provide additional qualitative results in \cref{fig:appendix_gen1} and \cref{fig:appendix_gen2}.
These results are uncurated samples of the classes listed as conditions.
These results (and \cref{fig:result_images}) are from our iMF-XL/2 model for 1-NFE ImageNet 256$\times$256 generation.
Following common practice, we present qualitative results using a CFG setting that favors the IS metric (emphasizing individual quality) at the expense of the FID metric (emphasizing diversity and distributional coverage); note that this tradeoff was impossible in the original MF, where the CFG is fixed.
Here, we set CFG as $\omega=6.0$ and CFG interval as $[t_\text{min}, t_\text{max}]=[0.2,0.8]$. This evaluation setting has an FID of 3.92 and an IS of 348.2.